% Generated draft: requires mathematical and source review.
\documentclass[11pt]{article}
\usepackage[margin=25mm]{geometry}
\usepackage{fontspec}
\setmainfont{Latin Modern Roman}
\usepackage{amsmath,amssymb,mathtools,unicode-math}
\setmathfont{Latin Modern Math}
\usepackage{enumitem,needspace,xurl,hyperref}
\setlength{\parindent}{0pt}
\setlength{\parskip}{5pt}
\setlength{\emergencystretch}{3em}
\newcommand{\N}{\mathbb{N}}
\newcommand{\Z}{\mathbb{Z}}
\newcommand{\Q}{\mathbb{Q}}
\newcommand{\R}{\mathbb{R}}
\newcommand{\C}{\mathbb{C}}
\newcommand{\F}{\mathbb{F}}
\newcommand{\A}{\mathbb{A}}
\newcommand{\PP}{\mathbb{P}}
\newcommand{\E}{\mathbb{E}}
\newcommand{\eps}{\varepsilon}
\newcommand{\dd}{\,\mathrm{d}}
\begin{document}
\section*{UnsolvedMath ID: 701761}
\textbf{Source problem code:} catalogue entry
\subsection{Definitions and mathematical statement}
For a group \(G\), define its derived series by
\[
G^{(0)}=G,\qquad G^{(k+1)}=[G^{(k)},G^{(k)}]\quad(k\geq 0),
\]
where \([H,H]\) is the subgroup generated by the commutators \(x^{-1}y^{-1}xy\) with \(x,y\in H\). The group \(G\) is soluble of derived length exactly \(n\) if
\[
G^{(n)}=\{1\}\quad\text{and}\quad G^{(n-1)}\neq\{1\}.
\]
It is torsion-free if
\[
\forall g\in G\;\forall m\in\Z_{\geq 1},\qquad g^m=1\Longrightarrow g=1.
\]
\textbf{Question (source formulation).} For every integer \(n\geq 3\), let \(G_\Gamma\) be a partially commutative soluble group of derived length \(n\), with defining graph \(\Gamma\), in the sense intended in Kourovka Notebook Problem 20.103. Must
\[
\forall g\in G_\Gamma\;\forall m\in\Z_{\geq 1},\qquad g^m=1\Longrightarrow g=1
\]
hold?

\textbf{Unresolved qualification.} The supplied statement says that the definition of \(G_\Gamma\) is similar to the derived-length-two case in Problem 17.104, but neither that definition nor the text of Problem 17.104 is supplied. Consequently, the class of groups quantified over is not fully specified by this record. In particular, no defining presentation or restriction on the defining graph, including finiteness, is imposed here without source support. The phrase ``derived length \(n\)'' is retained as exact derived length; identifying the intended construction requires consultation of the referenced definition.
\subsection{Short English statement}
Kourovka Notebook Problem 20.103, proposed by E. I. Timoshenko, asks whether every partially commutative soluble group of derived length at least three is torsion-free. The intended construction of these groups is referred to Problem 17.104 and is missing from the supplied record, so the precise class of groups requires clarification.
\subsection{Sources}
\begin{itemize}
\item Supplied local UnsolvedMath registry record: permanent ID \(701761\), exact source problem code \texttt{catalogue entry}, title \emph{Kourovka Notebook Problem 20.103}. The record labels the problem \texttt{open}; this is supplied registry metadata, not an independently verified assessment. The supplied record was updated at \texttt{2026-10-16T14:00:00Z}.
\item Attributed source: \emph{The Kourovka Notebook (Unsolved Problems in Group Theory)}, 20th issue (2022), Problem 20.103; also identified in the supplied background with the 21st issue edition (2026), \texttt{arXiv:1401.0300v48}. Proposed by E. I. Timoshenko in 2022. The supplied collection description names E. I. Khukhro and V. D. Mazurov as editors.
\item The statement refers to Kourovka Notebook Problem 17.104 for the derived-length-two definition. Its text is not included in the supplied record, and the referenced source has not been independently consulted here.
\end{itemize}
\end{document}
